\hypertarget{appendix-h-cpython-source-code-walkthrough-core-objects}{%
\chapter{Appendix H: CPython Source Code Walkthrough (Core
Objects)}\label{appendix-h-cpython-source-code-walkthrough-core-objects}}

This appendix provides a line-by-line analysis of the most critical C
functions in the CPython source code, allowing for an absolute
understanding of how the core data structures operate.

\hypertarget{h.1-objectslistobject.c-list_resize}{%
\subsection{\texorpdfstring{H.1 \texttt{Objects/listobject.c}:
\texttt{list\_resize}}{H.1 Objects/listobject.c: list\_resize}}\label{h.1-objectslistobject.c-list_resize}}

When you append to a list and it exceeds its current capacity, CPython
resizes the underlying array using an over-allocation strategy.

\begin{lstlisting}[language=C]
static int
list_resize(PyListObject *self, Py_ssize_t newsize)
{
    PyObject **items;
    size_t cur_allocated = (size_t)self->allocated;
    size_t allocated;

    if (cur_allocated >= (size_t)newsize && newsize >= (cur_allocated >> 1)) {
        assert(self->ob_item != NULL || newsize == 0);
        Py_SET_SIZE(self, newsize);
        return 0;
    }

    /* This over-allocation pattern is intended to give
       amortized O(1) performance for series of appends. */
    allocated = ((size_t)newsize + (newsize >> 3) + 6) & ~(size_t)3;
    if (newsize == 0)
        allocated = 0;

    items = (PyObject **)PyMem_Realloc(self->ob_item, allocated * sizeof(PyObject *));
    if (items == NULL) {
        PyErr_NoMemory();
        return -1;
    }
    self->ob_item = items;
    self->allocated = (Py_ssize_t)allocated;
    Py_SET_SIZE(self, newsize);
    return 0;
}
\end{lstlisting}

\begin{itemize}
\tightlist
\item
  \textbf{The Over-allocation Formula}:
  \passthrough{\lstinline!(newsize + (newsize >> 3) + 6) \& \~3!}. This
  ensures the list grows by about 12.5\% each time, plus a small
  constant, and remains aligned to a 4-item boundary.
\end{itemize}

\hypertarget{h.2-objectsdictobject.c-lookdict_unicode}{%
\subsection{\texorpdfstring{H.2 \texttt{Objects/dictobject.c}:
\texttt{lookdict\_unicode}}{H.2 Objects/dictobject.c: lookdict\_unicode}}\label{h.2-objectsdictobject.c-lookdict_unicode}}

This is the highly optimized lookup function for dictionaries where all
keys are Unicode strings (the most common case).

\begin{lstlisting}[language=C]
static Py_ssize_t
lookdict_unicode(PyDictObject *mp, PyObject *key, Py_hash_t hash)
{
    PyDictUnicodeEntry *ep0 = DK_UNICODE_ENTRIES(mp->ma_keys);
    size_t mask = DK_MASK(mp->ma_keys);
    size_t i = (size_t)hash & mask;
    PyDictUnicodeEntry *ep = &ep0[i];

    if (ep->me_key == NULL) return i;
    if (ep->me_key == key) return i;

    // ... Collision handling (linear probing with perturbation) ...
    for (size_t perturb = (size_t)hash; ; perturb >>= PERTURB_SHIFT) {
        i = (i << 2) + i + perturb + 1;
        ep = &ep0[i & mask];
        if (ep->me_key == NULL || ep->me_key == key) return i & mask;
    }
}
\end{lstlisting}

\begin{itemize}
\tightlist
\item
  \textbf{Optimization}: Note the use of
  \passthrough{\lstinline!(i << 2) + i!} (which is
  \passthrough{\lstinline!5*i!}). This is a fast way to perform the
  linear probing calculation without a slow multiplication instruction.
\end{itemize}

\begin{center}\rule{0.5\linewidth}{0.5pt}\end{center}
